\documentclass[11pt]{article}
\usepackage[margin=1in]{geometry}
\usepackage{enumitem}
% =============================================================================
% Preambles/header.tex — shared LaTeX/Beamer preamble for this project
%
% Use from a Beamer lecture by prepending:
%     \input{header}
% and compiling with TEXINPUTS=../Preambles:$TEXINPUTS (the /compile-latex
% skill does this automatically).
%
% PALETTE CONTRACT. Color names defined here MUST match the names in
% Quarto/theme-template.scss so Beamer and Quarto renderings use the same
% palette. The scripts/check-palette-sync.sh script enforces this.
%
% To customize: edit the HEX values below AND the matching $variable in
% Quarto/theme-template.scss, then run scripts/check-palette-sync.sh.
% =============================================================================

% -----------------------------------------------------------------------------
% Engine detection + encoding
%
% Our /compile-latex skill uses XeLaTeX, where inputenc is unnecessary and
% can produce encoding warnings. Only load inputenc under pdfTeX.
% -----------------------------------------------------------------------------
\usepackage{iftex}
\ifPDFTeX
  \usepackage[utf8]{inputenc}
\fi

\usepackage{amsmath, amssymb, amsthm}
\usepackage{booktabs}
\usepackage{graphicx}

% -----------------------------------------------------------------------------
% PALETTE — mirrors Quarto/theme-template.scss variable names.
% Load xcolor and define colors BEFORE anything that references them
% (notably hyperref, hypersetup, and the Beamer color assignments).
% -----------------------------------------------------------------------------
\usepackage{xcolor}

\definecolor{primary-blue}{HTML}{012169}      % headings, accents
\definecolor{primary-gold}{HTML}{B9975B}      % emphasis, borders
\definecolor{highlight-yellow}{HTML}{F2A900}  % markers, alerts
\definecolor{light-bg}{HTML}{E8EDF5}          % title / section backgrounds
\definecolor{jet}{HTML}{1A1A1A}               % body text

% Semantic colors (used in diagrams and annotations)
\definecolor{positive}{HTML}{15803D}          % good / observed / identified
\definecolor{negative}{HTML}{B91C1C}          % bad / problematic / bias
\definecolor{neutral}{HTML}{525252}           % reference / context

% Highlight accents (match .hi-* classes in the SCSS)
\definecolor{hi-slate}{HTML}{314F4F}
\definecolor{hi-green}{HTML}{2E8B57}
\definecolor{hi-red}{HTML}{C41E3A}

% Semantic aliases so skill templates and snippets can write
% \textcolor{accent}{...} without hard-coding "primary-blue".
\colorlet{accent}{primary-blue}
\colorlet{accent2}{primary-gold}

% -----------------------------------------------------------------------------
% Hyperref — loaded AFTER xcolor + palette so link colors resolve.
% -----------------------------------------------------------------------------
\usepackage{hyperref}
\hypersetup{colorlinks=true, linkcolor=primary-blue, urlcolor=primary-blue,
            citecolor=primary-blue, pdfborder={0 0 0}}

% -----------------------------------------------------------------------------
% Course code — set once per deck (after \input{header}, before \begin{document})
% via \coursecode{CS401}. Shown in the footer on every slide so a deck's course
% is identifiable without opening the title page. Empty by default.
% -----------------------------------------------------------------------------
\makeatletter
\newcommand{\coursecode}[1]{\gdef\@coursecodetext{#1}}
\gdef\@coursecodetext{}
\makeatother

% -----------------------------------------------------------------------------
% Beamer theme assignments (only apply under Beamer).
%
% We detect Beamer via \@ifclassloaded{beamer}, which is the documented,
% non-internal way. This must be wrapped in \makeatletter/\makeatother
% because \@ifclassloaded contains an @-catcode character.
% -----------------------------------------------------------------------------
\makeatletter
\@ifclassloaded{beamer}{%
  \usefonttheme{professionalfonts}
  \setbeamercolor{structure}{fg=primary-blue}
  \setbeamercolor{frametitle}{fg=primary-blue}
  \setbeamercolor{title}{fg=primary-blue}
  \setbeamercolor{subtitle}{fg=primary-gold}
  \setbeamercolor{author}{fg=primary-blue}
  \setbeamercolor{institute}{fg=neutral}
  \setbeamercolor{date}{fg=primary-gold}
  \setbeamercolor{itemize item}{fg=primary-blue}
  \setbeamercolor{itemize subitem}{fg=primary-gold}
  \setbeamercolor{alerted text}{fg=primary-gold}
  \setbeamercolor{block title}{fg=white, bg=primary-blue}
  \setbeamercolor{block body}{bg=light-bg}
  % Semantic block variants: block=definition/concept (blue), exampleblock=
  % worked example (green/positive), alertblock=key takeaway or caution (gold).
  % Keep to <=2 colored boxes per slide across ALL three variants (INV-7).
  \setbeamercolor{block title example}{fg=white, bg=positive}
  \setbeamercolor{block body example}{bg=light-bg}
  \setbeamercolor{block title alerted}{fg=white, bg=primary-gold}
  \setbeamercolor{block body alerted}{bg=light-bg}

  % Minimal footer: course code (left) + page X / N (right), no navigation clutter
  \setbeamertemplate{navigation symbols}{}
  \setbeamertemplate{footline}{%
    \color{neutral}\footnotesize\hspace{0.7em}\@coursecodetext\hfill
    \insertframenumber/\inserttotalframenumber\hspace{0.7em}\vspace{0.3em}}
}{}
\makeatother

% -----------------------------------------------------------------------------
% TikZ setup — libraries the snippet gallery uses
% -----------------------------------------------------------------------------
\usepackage{tikz}
\usetikzlibrary{arrows.meta, positioning, calc, decorations.pathreplacing,
                fit, shapes.geometric, backgrounds}

% Shared TikZ styles — snippets and hand-written diagrams can pick these up
% via `\begin{tikzpicture}[every node/.style={...}]` or by naming specific
% styles (e.g., `\node[dag-node] (X) at (0,0) {X};`).
\tikzset{
  % Node styles for causal DAGs and similar
  dag-node/.style = {
    draw=primary-blue, fill=light-bg, rounded corners=2pt,
    minimum width=1.2cm, minimum height=0.9cm, align=center, font=\small
  },
  decision-node/.style = {
    draw=primary-gold, fill=white, diamond, aspect=1.8,
    minimum width=1.8cm, minimum height=1.2cm, align=center, font=\small,
    inner sep=1pt
  },
  % Edge styles — observed vs counterfactual
  observed-edge/.style = {-{Latex[length=2.5mm]}, thick, draw=jet},
  counterfactual-edge/.style = {-{Latex[length=2.5mm]}, thick, draw=neutral, dashed},
  confound-edge/.style = {-{Latex[length=2.5mm]}, thick, draw=negative, dashed},
  % Data-point styles for plots
  observed-dot/.style = {fill=primary-blue, draw=primary-blue, circle, inner sep=1.2pt},
  counterfactual-dot/.style = {fill=white, draw=neutral, circle, inner sep=1.2pt}
}

% -----------------------------------------------------------------------------
% Convenience macros
% -----------------------------------------------------------------------------
\newcommand{\muted}[1]{\textcolor{neutral}{#1}}
\newcommand{\key}[1]{\textcolor{primary-gold}{\textbf{#1}}}
\newcommand{\good}[1]{\textcolor{positive}{#1}}
\newcommand{\bad}[1]{\textcolor{negative}{#1}}

% Standout frame macro: full-bleed dark background for section transitions.
% Beamer-only (uses \begin{frame}/\setbeamercolor/\usebeamerfont) -- do not
% call from an article-class file (e.g. Notes/<CODE>/*-notes.tex). \muted,
% \key, \good, \bad, and \coursecode above are plain xcolor/gdef and are
% safe to use from either class; verified when Notes/article-class support
% was added.
% Expand: \transitionslide{Your transition title}
\newcommand{\transitionslide}[1]{%
  \begingroup
  \setbeamercolor{background canvas}{bg=primary-blue}
  \begin{frame}[plain]
    \centering
    \vfill
    {\usebeamerfont{title}\color{white}\Large #1\par}
    \vfill
  \end{frame}
  \endgroup
}

% Section divider frame (Beamer-only), an alternative to \transitionslide with a
% darker two-line style: near-black (jet) background, a small white label line
% above a larger gold title, and a thin gold rule along the bottom edge.
% Expand: \sectiondivider{Section 2}{Pointers}
\newcommand{\sectiondivider}[2]{%
  \begingroup
  \setbeamercolor{background canvas}{bg=jet}
  \begin{frame}[plain]
    \vspace*{3.0cm}
    {\color{white}\ttfamily\large #1\par}
    \vspace{0.5cm}
    {\color{primary-gold}\LARGE #2\par}
    \begin{tikzpicture}[remember picture, overlay]
      \draw[primary-gold, line width=2.5pt]
        ($(current page.south west)+(0,0.1)$) -- ($(current page.south east)+(0,0.1)$);
    \end{tikzpicture}
  \end{frame}
  \endgroup
}

% -----------------------------------------------------------------------------
% Channel watermark — "Concepts That Click" (YouTube).
%
% Article-class files (CompetitiveExam Q/A sets, Notes, Assignments) get a
% light diagonal draft watermark on every page plus a centered footer naming
% the channel. Beamer decks get a subtle TikZ background watermark instead
% (the footline already carries the course code). Centralized here so every
% MCQ PDF is branded without per-file edits.
% -----------------------------------------------------------------------------
\makeatletter
\@ifclassloaded{article}{%
  \usepackage{draftwatermark}
  \SetWatermarkText{Concepts That Click}
  \SetWatermarkScale{1}
  \SetWatermarkFontSize{2cm}
  \SetWatermarkLightness{0.86}
  \SetWatermarkAngle{45}
  \usepackage{fancyhdr}
  \pagestyle{fancy}
  \fancyhf{}
  \fancyfoot[C]{\footnotesize\textcolor{neutral}{Concepts That Click~|~YouTube: \href{https://www.youtube.com/@concepts.thatclick}{@concepts.thatclick}}}
  \renewcommand{\headrulewidth}{0pt}
  \renewcommand{\footrulewidth}{0pt}
  \fancypagestyle{plain}{%
    \fancyhf{}%
    \fancyfoot[C]{\footnotesize\textcolor{neutral}{Concepts That Click~|~YouTube: \href{https://www.youtube.com/@concepts.thatclick}{@concepts.thatclick}}}%
    \renewcommand{\headrulewidth}{0pt}%
    \renewcommand{\footrulewidth}{0pt}%
  }
}{}
\@ifclassloaded{beamer}{%
  \setbeamertemplate{background}{%
    \begin{tikzpicture}[remember picture, overlay]
      \node[rotate=45, opacity=0.055, align=center]
        at (current page.center)
        {\Huge\textcolor{neutral}{Concepts That Click}};
    \end{tikzpicture}%
  }
}{}
\makeatother 
\coursecode{JKSSB}
\renewcommand{\thesection}{43.\arabic{section}}
\newtheorem{example}{Example}[section]
\newenvironment{definitionbox}[1]{%
  \par\noindent\textbf{Definition (#1).}\ \itshape
}{\par\vspace{0.3em}}
\title{Hacking, Phishing and Spoofing --- Detailed Lecture Notes}
\author{JKSSB Computer Awareness}
\date{\today}

\begin{document}
\maketitle

\section{Why attacks succeed}
Most attacks on ordinary computer users do not defeat the machine by
force. They persuade the person in front of the screen to hand over the
key: a password, a One-Time Password, a click, or a file download. The
attacker needs a weakness to use and a moment in which the victim acts
without verifying. The weakness may be technical, such as a short reused
password or outdated software, or human, such as trust in an
authoritative-sounding caller, hurry, or fear of an account being
blocked. This lesson stays at a safe awareness level: it names each
attack, shows how to recognise it, and teaches the correct defensive
response. It does not teach how to carry out any attack.

\begin{definitionbox}{Attack chain}
A \textbf{vulnerability} is a weakness in software, settings, or human
habit that can be misused. An \textbf{exploit} is a method that takes
advantage of such a weakness. An \textbf{attack} is the complete harmful
act against a person or system. The \textbf{payload} is the harmful
element the attack delivers, such as stolen credentials, stolen money,
or a malicious file.
\end{definitionbox}

\begin{center}
\begin{tabular}{p{0.24\textwidth}p{0.66\textwidth}}
\toprule
\textbf{Term} & \textbf{Meaning in this lesson} \\
\midrule
Hacking & Gaining unauthorised access to a computer, account, or network. \\
Phishing & A fake message, mail, or call imitating a trusted sender to steal credentials or money. \\
Spoofing & Faking an identity, such as a sender address, caller number, or website address. \\
Social engineering & Tricking a person, rather than breaking software, into giving access. \\
Password attack & An attempt to discover a password by guessing, a word list, exhaustive trial, or observation. \\
Malicious link & A web address that leads to a fake page or an unwanted download. \\
Malicious attachment & A file that harms the computer or steals data when opened. \\
OTP & One-Time Password: a single-use code; never shared with anyone. \\
\bottomrule
\end{tabular}
\end{center}

The rest of this lesson uses these terms in one consistent sense.
\textbf{Hacking} is the broad act of unauthorised access.
\textbf{Phishing} is fraud by imitation that asks the victim to act.
\textbf{Spoofing} is the forged identity that makes the fraud look
genuine. \textbf{Social engineering} is the human manipulation behind
both. A \textbf{vulnerability} is the weakness, the \textbf{exploit} is
the method, and the \textbf{payload} is the harm delivered.

\section{Hacking}
\textbf{Hacking} means gaining \textbf{unauthorised access} to a
computer system, an online account, or a network. A \textbf{hacker} is
the person who performs this act. Access without permission is unlawful
even when nothing is copied or damaged; the absence of permission is
what defines the act, not the amount of harm that follows.

An \textbf{ethical hacker} is a security tester who probes systems
\textbf{with written permission} from the owner in order to find and fix
weaknesses before criminals exploit them. The permission is the entire
difference: the same technical step performed with written authorisation
is a security test, and performed without it is an attack.

\begin{definitionbox}{Hacking}
Gaining access to a computer, account, or network without the owner's
permission.
\end{definitionbox}

\begin{example}
A student opens a colleague's mailbox by guessing the password and
reads the mail. Even though nothing is deleted, this is hacking,
because the access was unauthorised. A bank that hires a tester to try
its login page, with a signed agreement, is instead conducting an
ethical security test.
\end{example}

\section{Password attacks}
Passwords remain the most common lock on an account, so attackers try
to discover them first. Four methods appear repeatedly in examinations.

\textbf{Guessing} tries obvious choices: the user's name, birth year,
phone number, or common strings such as \texttt{password123}.
\textbf{Dictionary attack} tries every entry in a list of common
passwords automatically, which succeeds quickly against short familiar
words. \textbf{Brute-force attack} tries every possible combination of
characters until one works; it always succeeds eventually, but a longer
password makes the search impractically slow. \textbf{Shoulder surfing}
is direct observation: watching the keyboard or screen while the victim
types a password or PIN.

The defence follows from the methods. Use a \textbf{long} password that
mixes letters, digits, and symbols, and never reuse one password across
accounts. Never disclose a password or an \textbf{OTP (One-Time
Password)} to anyone, including a caller who claims to represent the
bank or the office. Lock the screen when leaving the desk, shield the
keypad while typing, and switch on \textbf{two-step verification} (also
called two-factor authentication) wherever it is offered, so that the
password alone is not enough to enter the account.

\begin{definitionbox}{Brute-force vs dictionary attack}
A \textbf{dictionary attack} tries each word in a ready list of common
passwords. A \textbf{brute-force attack} tries every possible character
combination in turn.
\end{definitionbox}

\section{Social engineering}
\textbf{Social engineering} is the practice of tricking a person into
granting access, instead of breaking the software technically. Two
named forms recur. \textbf{Pretexting} invents a false identity and
story --- a bank officer verifying records, an IT support engineer
fixing the account --- to extract a password, OTP, or remote-access
permission. \textbf{Baiting} offers a lure --- a prize, a free recharge,
free software --- that triggers the victim to click a link or download
a file.

The standard signal is manufactured urgency: words such as
\textbf{urgent}, \textbf{blocked}, \textbf{suspended}, and \textbf{verify
now}, combined with a demand to act immediately without checking. The
safe response is fixed: slow down, verify the claim through an
independently known channel (the published helpline or the real
website), and never disclose a password or OTP under pressure.

\begin{definitionbox}{Social engineering}
Manipulating a person, rather than a machine, into revealing
credentials or granting access.
\end{definitionbox}

\section{Phishing}
\textbf{Phishing} is a fake message, email, text, or call that
\textbf{imitates a trusted sender} --- a bank, an office, a delivery
service --- in order to steal passwords, OTPs, card details, or money.
The message always contains a call to action: click the enclosed link,
open the attached file, or call back the given number. The link leads
to a look-alike login page that records whatever is typed; the
attachment installs the payload when opened.

Named channels help classification. Fraudulent text messages are called
\textbf{smishing} (SMS phishing) and fraudulent voice calls are called
\textbf{vishing} (voice phishing), but both are phishing by the same
definition: imitation plus a request for action. The recognition signs
are stable: an urgent or threatening tone, spelling and grammar errors,
a sender address that does not match the claimed organisation, and a
link whose real address differs from the genuine website.

\begin{definitionbox}{Phishing}
A fraudulent message that imitates a trusted sender and asks the victim
to click, open, reply, or pay, in order to steal credentials or money.
\end{definitionbox}

\begin{example}
A clerk receives a mail headed ``Salary slip ready --- download within
24 hours or the account will be blocked.'' The sender shows the office
name, but the reply address is an unknown mailbox and the enclosed
link points to a misspelled portal. The mail is phishing: imitation (the
office name), pressure (blocking within 24 hours), and a call to action
(the link). The safe response is to ignore the link, open the real
portal by typing its known address, and report the message.
\end{example}

\section{Spoofing}
\textbf{Spoofing} means \textbf{faking an identity} so that a message,
call, or website appears to come from someone it does not. Three forms
cover the examination scope. \textbf{Email spoofing} forges the
``From'' address so the mail displays a trusted name that never sent
it. \textbf{Caller ID spoofing} forges the number shown on the phone so
the call appears to come from a bank or office. \textbf{Website (URL)
spoofing} registers a look-alike address --- for instance
\texttt{bank-verlfy.com} in place of the genuine
\texttt{bank-verify.com} --- and hosts a fake login page there.

Spoofing is best understood as a disguise rather than a complete fraud.
A phishing message very often \textbf{uses} spoofing --- a forged
sender address --- to look trustworthy, but spoofing can also appear
without any message at all, as in a forged website waiting for
visitors. The check is the same in every case: confirm the identity
through a known channel instead of trusting what is displayed.

\begin{definitionbox}{Spoofing}
Presenting a forged sender address, caller number, or website address
as though it were genuine.
\end{definitionbox}

\section{Phishing vs spoofing vs hacking}
These three terms form the mandatory near-neighbour set of the lesson,
and examination options deliberately mix them. \textbf{Hacking} is the
broadest: any unauthorised access. \textbf{Phishing} is fraud by
imitation that asks the victim to act: click, open, reply, or pay.
\textbf{Spoofing} is the forged identity itself: the faked address,
number, or site. A single incident can contain all three --- a forged
sender address (spoofing) carries a fake bank mail asking for the OTP
(phishing), and the stolen OTP is then used to enter the account
(hacking) --- but each word answers a different question: was access
taken without permission, was a fake message asking for action, or was
an identity forged?

\section{Malicious links and attachments}
A \textbf{malicious link} is a web address that leads to a fake login
page or starts an unwanted download. Before clicking, hover over the
link with the mouse to reveal the real address, or long-press it on a
phone to preview the destination. Safer still, ignore the link entirely
and type the website's known address by hand. Shortened addresses,
misspelled brand names, and extra words inserted into a familiar
address, arriving inside an unexpected or urgent message, are warning
signs.

A \textbf{malicious attachment} is a file that harms the computer or
steals data once opened. Never open an unexpected attachment, even when
the displayed sender is a known name, because that name may itself be
spoofed; confirm with the sender through a separate channel first. Be
especially cautious with documents that ask the reader to ``enable
macros'' or ``enable content'' on opening, since that step is what
activates many payloads. In short: the link is the address, the
attachment is the file, and both can carry the payload.

\section{Safe habits and response after targeting}
Four habits block the large majority of these attacks. First, verify
every urgent request through a \textbf{known channel} --- the published
helpline or the manually typed website --- before acting. Second, keep
the operating system, browser, and \textbf{antivirus} software updated,
and scan external drives before opening their files. Third, keep a
regular \textbf{backup} of important files so that one incident cannot
destroy the only copy. Fourth, use the mail client's \textbf{report
phishing} option on a suspicious message rather than replying to it or
forwarding it to colleagues.

If a suspicious page has already been opened, do not type any password,
OTP, or card details into it. Close the page, delete the message, and
do not call back any number printed inside it. Change the affected
password by visiting the real website directly, inform the bank or
office through its published helpline, and preserve a record --- a
screenshot and the sender's address --- for the complaint. The order is
fixed: stop, verify on the real channel, change credentials, report.

\section{Exam retrieval and common confusions}
Six distinctions prevent most errors on this topic.
\begin{enumerate}
  \item Phishing \textbf{imitates} a trusted sender to steal data or
    money; spoofing \textbf{forges} the identity it wears; hacking is
    the broader \textbf{unauthorised access}.
  \item OTP stands for \textbf{One-Time Password}. No genuine bank,
    office, or support desk ever asks for it by call or message.
  \item A \textbf{dictionary} attack tries a ready word list; a
    \textbf{brute-force} attack tries every combination.
  \item A \textbf{vulnerability} is the weakness, an \textbf{exploit} is
    the method that uses it, and the \textbf{payload} is the harm
    delivered.
  \item Ethical hacking requires \textbf{written permission}; without
    it, access is unauthorised.
  \item The safe response to any urgent message is fixed: do not click,
    verify through a known channel, and report.
\end{enumerate}

\begin{example}
An item describes a call in which the displayed number matches the
bank's helpline and the caller demands the OTP to ``unblock'' the
account. Trace the definitions: the displayed number is forged
(spoofing), the call imitates the bank and demands the OTP (phishing by
vishing, using social engineering), and any resulting entry into the
account is unauthorised access (hacking). The correct answer names
phishing, and the correct action is to refuse, disconnect, and verify
through the published helpline.
\end{example}

\subsection*{Exercises}
\begin{enumerate}[label=\arabic*.]
  \item Define hacking, and state what separates an ethical security
    test from an attack.
  \item Distinguish a dictionary attack from a brute-force attack.
  \item Define phishing and list four recognition signs of a phishing
    message.
  \item Distinguish phishing from spoofing in one sentence each, and
    explain how they combine.
  \item State the safe response when an unexpected message demands an
    OTP urgently.
\end{enumerate}

\subsection*{Solutions}
\begingroup\small
\begin{enumerate}[label=\arabic*.]
  \item Hacking is gaining unauthorised access to a computer, account,
    or network. An ethical security test is performed with the owner's
    written permission in order to find and fix weaknesses; without
    that permission the same act is an attack.
  \item A dictionary attack tries each entry of a ready list of common
    passwords; a brute-force attack tries every possible character
    combination in turn, so longer passwords resist it.
  \item Phishing is a fraudulent message imitating a trusted sender in
    order to steal credentials or money. Signs: urgent or threatening
    tone, spelling errors, a sender address that does not match the
    claimed organisation, and a link whose real address differs from
    the genuine site.
  \item Phishing is fraud by imitation that asks the victim to act;
    spoofing is the forged identity itself. A phishing message often
    wears a spoofed sender address to look trustworthy.
  \item Do not click, reply, or disclose the OTP; verify through an
    independently known channel (published helpline or manually typed
    website), change credentials on the real site if exposed, and
    report the message.
\end{enumerate}
\endgroup
\end{document}
